\documentclass[10pt,a4paper,twocolumn]{article}

\usepackage[T1]{fontenc}
\usepackage[utf8]{inputenc}
\usepackage{newtxtext,newtxmath}
\usepackage[a4paper,top=18mm,bottom=19mm,left=17mm,right=17mm,columnsep=6mm]{geometry}
\usepackage{microtype}
\usepackage{booktabs}
\usepackage{tabularx}
\usepackage{enumitem}
\usepackage{titlesec}
\usepackage{xcolor}
\usepackage{hyperref}
\usepackage{balance}

\definecolor{duckgreen}{HTML}{0B5141}
\hypersetup{colorlinks=true,linkcolor=duckgreen,urlcolor=duckgreen,citecolor=duckgreen,
pdfauthor={Tair Kezdekbayev, Alvaro Mendez Li, Tan Kayra Erol}}
\setlength{\parindent}{1em}
\setlength{\parskip}{0pt}
\setlength{\columnsep}{6mm}
\setlist{nosep,leftmargin=*,topsep=2pt}
\titleformat{\section}{\normalfont\scshape\large}{\thesection}{0.55em}{}
\titleformat{\subsection}{\normalfont\itshape\normalsize}{\thesubsection}{0.55em}{}
\titlespacing*{\section}{0pt}{8pt}{3pt}
\titlespacing*{\subsection}{0pt}{5pt}{2pt}

\newcommand{\code}[1]{\texttt{#1}}
\newcommand{\result}[1]{\textbf{#1}}

\begin{document}

\twocolumn[
\begin{@twocolumnfalse}
\begin{center}
    {\LARGE\bfseries Duckietown -- Vision-Guided Navigation and Live Chat\par}
    \vspace{7pt}
    {\large Tair Kezdekbayev \quad Alvaro Mendez Li \quad Tan Kayra Erol\par}
    \vspace{2pt}
    {\itshape Department of Computer Engineering\par}
    {\itshape Technical University of Munich\par}
    {Heilbronn, Germany\par}
\end{center}
\vspace{7pt}
\end{@twocolumnfalse}
]

\section{Abstract}

This project implements autonomous right-lane navigation for a Duckiebot using
ROS~1, monocular camera perception, a directed road map, and a desktop companion
application. The onboard controller segments yellow, white, and red road
markings, estimates lane position, regulates differential wheel commands, and
switches to bounded junction manoeuvres when lane borders disappear. A directed
A* planner calculates the shortest legal route between red-line approaches.
The Windows application combines route selection, live camera diagnostics,
status reporting, a working local live chatbot, and a direct stop control. The software was matched to the
robot's installed ROS Noetic environment. Straight driving, a left curve-to-straight
transition, a sharp right road bend, and individual straight, left, and right
junction traversals were demonstrated. The chatbot provides immediate timed pauses, straight-road speed profiles and
validated junction instructions. Its command rules run locally without model
training or a cloud service.

\section{Introduction}

Duckietown provides a compact autonomous-driving platform in which perception,
control and distributed software can be studied together
\cite{paull2017duckietown}. The objective of this project was to make one robot,
\emph{duck2}, follow the right lane and reach a selected red-line destination on
a known course. The solution had to operate with the camera and wheel interfaces
already installed on the robot, tolerate missing borders inside intersections,
and expose an understandable interface to the operator.

The resulting system divides responsibility between the robot and laptop. The
robot performs all time-critical perception, steering, wheel output, and safety
checks. The laptop selects a directed start, computes a route, sends validated
high-level commands, maintains a live-chat session, and visualises camera and controller data. Communication
uses an authenticated SSH tunnel, so the HTTP command and camera services remain
bound to loopback on the robot.

\section{System Architecture}

The robot host runs Ubuntu 18.04.6 on ARM64. Its installed ROS container
provides ROS Noetic, Python 3.8.10, OpenCV 4.2.0, and NumPy 1.17.4. Development images use the digest-pinned
\code{duckietown/dt-ros-commons:v4.3.0} base for both AMD64 and ARM64. The
camera input is \code{/duck2/camera\_node/image/compressed} with
\code{sensor\_msgs/CompressedImage}; wheel requests use
\code{/duck2/wheels\_driver\_node/wheels\_cmd} with
\code{duckietown\_msgs/WheelsCmdStamped}.

The project consists of four main parts:

\begin{enumerate}
    \item a ROS camera and lane-processing node;
    \item a navigation state machine for road segments and junctions;
    \item a command gateway, camera gateway, and independent watchdog; and
    \item a Windows companion application for map planning and supervision.
\end{enumerate}

Only the lane follower publishes wheel requests during project operation. The
gateway accepts structured high-level commands, while the companion subscribes
to status and camera data without becoming a second wheel controller.

\section{Lane Perception}

\subsection{Colour segmentation}

Each compressed frame is decoded in OpenCV, converted from BGR to HSV, and
cropped to a lower road region. The deployed colour intervals are

\begin{center}
\begin{tabular}{lll}
\toprule
Marking & Lower HSV & Upper HSV \\
\midrule
Yellow & $(20,70,80)$ & $(35,255,255)$ \\
White  & $(0,0,150)$ & $(180,55,255)$ \\
\bottomrule
\end{tabular}
\end{center}

Connected components and row-wise boundary samples remove small fragments and
separate longitudinal lane edges from transverse markings. The controller uses
ordered yellow-left and white-right evidence whenever both are available. A
short temporal width estimate supports partial observations; fresh frames with
insufficient road evidence still remain distinct from missing camera data.

\subsection{Lane geometry}

Let $x_y(r)$ and $x_w(r)$ be the detected yellow and white boundary positions at
image row $r$, and let $W$ be image width. The normalised lane centre is

\begin{equation}
 c(r)=\frac{x_y(r)+x_w(r)}{2W}, \qquad e(r)=c(r)-c_{\mathrm{ref}} .
\end{equation}

Matched image rows in the junction corridor estimator prevent perspective
from mixing unrelated depths.
The change in $c(r)$ between near and far rows gives an image-based heading
estimate, while the near value gives lateral displacement. The ordinary lane reference is $c_{\mathrm{ref}}=0.441$; the dedicated
straight-junction approach uses a separately tuned reference of $0.49$. These
image references are controller parameters rather than metric lane coordinates.

\subsection{Steering control}

Lane error is low-pass filtered,

\begin{equation}
 \hat e_t=\alpha e_t+(1-\alpha)\hat e_{t-1}, \qquad \alpha=0.20,
\end{equation}

then mapped to bounded steering with gain scheduling, a deadband, and slew-rate
limiting. For centre speed $v$ and steering request $\delta$, the wheel command
pair is

\begin{equation}
 u_L=v-\delta, \qquad u_R=v+\delta .
\end{equation}

The continuous launcher uses $v=0.09$, a wheel cap of $0.20$, a steering cap of
$0.11$, and proportional gains of $0.35$ near centre and $0.75$ farther away.
These values are normalised driver requests rather than metric velocities. The
installed driver adds a minimum PWM offset; consequently, small numerical wheel
differences do not always produce equally visible changes in motion.

\section{Road and Junction Navigation}

\subsection{Smooth and sharp roads}

Smooth curves remain under continuous visual lane control. A confirmed left
road curve can request the supported $0.03/0.20$ left/right profile; this boost
is gated by curve evidence and is not applied to every straight-road frame. A sharp right road
bend is recognised only during ordinary following, after recent valid lane
evidence and a confirmed disappearance pattern at the right boundary. The
bounded corner state first preserves forward progress and then applies the
physically verified $0.20/0.00$ left/right pivot request. It returns to visual
centering after a stable outgoing corridor appears. Corner recognition is
disabled while a junction manoeuvre is active, preventing both state machines
from controlling the same event.

\subsection{Red-line stopping}

A red transverse component near the bottom of the camera region latches a stop.
The robot outputs zero for a two-second dwell and waits for the route controller
to authorise one valid exit. A red line is therefore a navigation event rather
than merely a colour cue. Persistent images of the same departure line cannot
advance the route twice.

\subsection{Unmarked intersections}

The course has areas inside intersections where neither longitudinal boundary
is visible. These areas are handled as explicit states:

\begin{enumerate}
    \item \emph{approach}: ordinary visual lane following;
    \item \emph{stopped}: latched zero output and route validation;
    \item \emph{crossing}: a direction-specific bounded motion profile;
    \item \emph{reacquisition}: search for an ordered outgoing corridor; and
    \item \emph{following}: visual control resumes after stable alignment.
\end{enumerate}

Missing lane markings are permitted only in the authorised crossing and search
states. Straight crossing uses a stronger equal-wheel request with limited
encoder balancing. The left profile reproduces the verified gradual curve with
an outer-wheel advantage. The right profile uses forward clearance followed by
the calibrated sharp pivot. Reacquisition requires near-field support, correct
yellow/white order, acceptable lateral and heading errors, and a stability
interval. Every manoeuvre has a deadline.

\section{Route Planning}

The supplied course is represented as five junctions, $A$--$E$, with directed
approaches such as $A\!\rightarrow\!B$. A state stores the previous and current
junction, which fixes the robot's heading and right-lane position. Legal forward transitions
are labelled \emph{left}, \emph{straight}, or
\emph{right}. Red-line approaches, rather than undirected junction points, are
used as selectable destinations.

The laptop applies A* search \cite{hart1968astar}. Every crossed junction has
unit cost, and a relaxed graph distance is the admissible heuristic. Therefore,
the result minimises the number of junctions while respecting directed lanes.
Before execution, the route is converted into a sequence of turn instructions.
The next instruction is released only for the current run and junction index;
route progress advances only after the robot reports the corresponding outgoing
lane.

\section{Companion Application}

The \code{duck2\_companion.py} desktop application combines two work areas. The
map view lets the user choose a directed starting lane and destination red line,
shows the A* route, and requires explicit placement confirmation before Start.
The camera/status view displays the normal image, colour mask, or diagnostic
overlay together with controller phase, next junction, wheel request, and stop
reason. A prominent Stop control remains available throughout a run.

The application connects to loopback ports 8765 and 8766 through a strict-key
Windows OpenSSH tunnel. The robot-side services also bind to loopback. This
keeps command transport separate from ROS network discovery and makes session
startup repeatable after laptop or robot reboots.


\section{Live Chatbot}

\subsection{Implementation and command flow}

The live chatbot is integrated into the companion's camera/status panel. The
normal application enables it by default. The
operator types a message and presses Send or Enter; the input is not populated
with a scripted demonstration. \code{live\_chat.py} normalises text, recognises
an allowed intent and applies it to the current session. The UI reports actual
controller acknowledgments and state changes. It does not claim that an
accepted wheel request proves physical movement.

The laptop keeps the future turn queue and validates proposed changes against
the same directed map used by A*. A request such as ``go straight at the next
junction'' replaces the remaining queue with one legal straight instruction.
It changes the decision at a red-line stop; it does not disable lane steering
on subsequent curves. A turn already accepted for a crossing is completed
before another queued turn can be dispatched.

At each red stop the robot observes its two-second dwell and accepts one
instruction tied to the current run, approach and junction index. Repeated
delivery uses the same identifier, preventing a lost acknowledgment from
consuming two turns. The laptop updates its reported lane only after matching
the outgoing-lane completion report. With no queued instruction, the chatbot
asks where to go next and displays a 60-second waiting limit measured from
arrival. Invalid messages do not restart that waiting clock.

\subsection{Supported commands}

\begin{table}[h]
\centering
\caption{Examples of the implemented live command grammar.}
\label{tab:chat}
\footnotesize
\begin{tabularx}{\columnwidth}{@{}p{0.44\columnwidth}X@{}}
\toprule
Message & Effect \\
\midrule
\code{stop for 7s} & Pause immediately on acceptance, then resume after seven seconds if healthy. \\
\code{go straight at the next junction} & Replace pending turns with straight. \\
\code{go left}, \code{go right} & Choose the next legal exit. \\
\code{left after that} & Append left to the future queue. \\
\code{clear turns} & Empty the queue; ask at the next red stop. \\
\code{slow speed}, \code{normal speed}, \code{fast speed} & Select a straight-road speed profile. \\
\code{speed up}, \code{slow down} & Move one profile higher or lower. \\
\code{status} & Show reported position and pending instructions. \\
\code{stop}, \code{pause} & Pause awaiting Continue for at most 30 seconds. \\
\code{continue}, \code{resume} & Resume a pause without bypassing a red stop. \\
\code{quit}, \code{end run} & End the session and cancel automatic resumption. \\
\bottomrule
\end{tabularx}
\end{table}

The slow, normal and fast profiles request centre speeds of $0.09$, $0.10$ and
$0.11$, respectively, with a straight-profile wheel cap of $0.12$. Normal is
the initial selection. These profiles apply only when stable visual evidence
classifies the road as straight; curve and junction profiles keep their
separate settings. Values are normalised driver inputs, not measured speeds.

A timed pause requests zero output in the robot's command callback before
acknowledgment, without waiting for another camera frame. Its duration is
measured with the robot's monotonic clock. Network delivery and mechanical
braking still introduce latency. The queue survives the pause, while camera,
heartbeat and watchdog checks remain active. A fault or Stop ends the run and
prevents automatic resumption. Timed requests accept positive durations up to
3,600 seconds; the shorter 30-second limit applies to an untimed pause.

\subsection{How the chatbot was developed}

The chatbot uses a hand-written, deterministic command grammar. It was
developed by specifying intents and writing phrase-recognition rules for the
supported messages and current route state. It has no trained
neural network, fitted weights, fine-tuning dataset or online learning stage.
Command rules were refined with examples rather than machine-learning training.

Text processing tolerates capitalisation, common punctuation and supported
polite phrasing. Explicit route state supplies context for ``after that'',
relative speed changes and the next junction. Unsupported wording and
ambiguous requests produce clarification; arbitrary spelling mistakes are not
guessed. No API key or downloaded language model is needed.

\section{Safety Supervision}

Movement requires all of the following: a fresh camera frame, a healthy
controller, sole wheel-publisher ownership, a current laptop heartbeat, valid
route state, and explicit user confirmation. The independent watchdog monitors
camera and status freshness, wheel feedback, encoders, publisher ownership, and
parent-process health. The controller requests zero after a manual Stop, stale data or an expired
manoeuvre. On a detected fault, the independent watchdog asserts the installed
driver emergency stop and ends the temporary controller process group.

Cleanup waits for temporary publishers to exit and confirms zero wheel
feedback. Restoration of normal robot control is a deliberate step after the
project session ends; it is not an automatic app reconnection action.

\balance
\section{Operational Limitations}

The software remains map-relative rather than globally localised. The operator
must select the actual starting approach. Physically relocating the robot
requires a new session rather than relying on its previous reported position.
The normal application uses the selected map route as its initial queue and
permits live overrides at subsequent red stops. Individual maneuvers have been
demonstrated, but full-route reliability remains incomplete: insufficient curve
following and heading errors can still interrupt navigation. Implemented command
delivery does not establish reliable containment under every physical condition.

\section{Reproducible Operation}

Windows OpenSSH uses a protected key and strict host verification. After boot,
run the connection check, then \code{tools/Start-Duck2-DrivingMode.cmd} while
stationary. Preparation stages current source, establishes exclusive wheel
ownership and leaves the controller stopped. Open
\code{laptop/Start-Duck2Companion.cmd}, connect, select the actual starting lane
and destination, confirm placement and press Start. The operator types each
chat message; the app does not prefill instructions. Stop ends the run.

Source changes require a newly prepared container and reopened app. Camera-only
mode uses a driving-disabled backend. Docker Desktop is needed for optional
local image builds, not source-mounted operation. Repository guides provide
installation, shutdown and connection-recovery instructions.

\section{Conclusion}

The project combines camera perception, lane steering, red stops, junction
traversal, directed A* planning and local live chat. Validated text commands
pause the robot, select straight-road speed profiles and revise future junction
choices. Language interpretation remains outside the wheel-control loop.
The implemented software and demonstrated component maneuvers remain subject
to the complete-route perception and alignment limitations described above.

\begin{thebibliography}{9}
\bibitem{paull2017duckietown}
L. Paull \emph{et al.}, ``Duckietown: An open, inexpensive and flexible platform
for autonomy education and research,'' in \emph{Proc. IEEE Int. Conf. Robotics
and Automation}, 2017, pp. 1497--1504.

\bibitem{quigley2009ros}
M. Quigley \emph{et al.}, ``ROS: An open-source Robot Operating System,'' in
\emph{ICRA Workshop on Open Source Software}, 2009.

\bibitem{hart1968astar}
P. E. Hart, N. J. Nilsson, and B. Raphael, ``A formal basis for the heuristic
determination of minimum cost paths,'' \emph{IEEE Trans. Systems Science and
Cybernetics}, vol. 4, no. 2, pp. 100--107, 1968.

\bibitem{opencv}
G. Bradski, ``The OpenCV Library,'' \emph{Dr. Dobb's Journal of Software Tools},
2000.

\bibitem{projectrepo}
T. Kezdekbayev, A. Mendez Li, and T. K. Erol, ``AlTaTa Duckietown Final,''
software repository, 2026.
\href{https://github.com/DuckietownTUM/AlTaTa-Duckietown-Final}{GitHub: DuckietownTUM/AlTaTa-Duckietown-Final}.
\end{thebibliography}

\end{document}
